\chapter{The Comtet Majumdar family and the uniform law}\label{chap:cg}\chaptermark{The Comtet Majumdar family}

\providecommand{\CGLi}{\operatorname{Li}}
\providecommand{\CGLP}{\mathcal{L}\mathcal{P}}

Comtet and Majumdar \cite{ComtetMajumdar05} proved, in an appendix on the maximum of a random walker, an identity between an integral of the logarithm of $(1-e^{-k^\mu})/k^\mu$ and the value $\zeta(1/\mu)$ of the Riemann zeta function, and they proved it through a cutoff regularisation. This chapter asks what that identity and that regularisation can do for the programme of the book. The programme is to approximate $\Xi$ by entire functions whose reciprocals are already Laplace transforms of generalised gamma convolutions (GGC), and then to pass to the limit. The answer has two halves. On the arithmetic side the answer is negative: the shifted form of the regularisation of Comtet and Majumdar, and every decreasing damping of the coefficients of the theta series, produces functions that are not in the Laguerre P\'olya class, for structural reasons that are proved below. On the Thorin side the same shift, read as a tilt of a L\'evy measure, gives a positive object: for the uniform law the Thorin density of the tilted law is strictly positive for every positive tilt, and its Mellin transform is the regularised zeta function of the family. The uniform law is therefore a solved model of the question that the book asks for $\xi$ in Chapter~\ref{chap:basepoint}, where the corresponding positivity is proved for every basepoint $\alpha\ge1$.

The chapter develops a source note by Nick Polson \cite{P26cg}, together with the computations done for this book. Its numbers were recomputed here by a script written from the mathematics alone (the scripts of the source were not used; the script of this chapter is \texttt{code/cg\_checks.py}, with its output in \texttt{data/cg\_checks\_*.json} and the logs beside them, listed in \texttt{code/README.md}). Every number printed below is one of those recomputations unless it is marked unchecked, and where the recomputation differed from the source the recomputed value is printed. The tables of zeros of the damped kernels (Section~\ref{cg:sec:kernel}) were not recomputed and are marked unchecked. The notation is the book's, with the following warnings. The kernel $\Phi$ is the kernel of \eqref{eq:ch2:Phi}, so that $\Xi(t)=\int_\R\Phi(u)\cos(tu)\,du$, and the sum over $n$ inside $\Phi$ carries the terms $\Phi_n$ of Section~\ref{cg:sec:kernel}. The damping parameter is $\delta$ (elsewhere in the book $\delta$ may denote an off line displacement, which does not occur here). The Poisson kernel of the book is $P_a$ with the factor $1/\pi$ (Chapter~\ref{chap:evaluation}); the unnormalised Lorentzian of this chapter is $\ell_\varepsilon(x)=\varepsilon/(\varepsilon^2+x^2)=\pi P_\varepsilon(x)$. The Thorin density of the uniform law is $\kappa_\delta$; for $\xi$ the book's $D_\alpha$ of Chapter~\ref{chap:polyazeros} plays the role of $\kappa^\xi_\alpha$. The constant of Comtet and Majumdar is written $C_\mu$ (the letter $\gamma$ is kept for ordinates of zeros), the exponential variables are $e_k$ as in Chapter~\ref{chap:evaluation}, and the integer part is written $\lfloor t\rfloor$ (the letter $N$ is kept for zero counts). The de Bruijn Newman flow is written with the time $t$ and the kernel $\Phi_{\mathrm{dBN}}(u)=\frac12\Phi(2u)$ of Section~\ref{ch:dbn}, not with the parameter $\lambda$ of the source, whose normalisation differs from the literature (Section~\ref{cg:sec:dbn}).

Several parts of the source are already in the book and are cited, not restated. The Thorin form of the reciprocal $\xi$ function with a basepoint, its L\'evy form, the compensated density $v_\alpha-b_\alpha$ and its positivity for every $\alpha\ge1$ (Proposition~\ref{bp:prop:levy}, Proposition~\ref{bp:prop:valpha}, Theorem~\ref{thm:dbeta:minall}, Remark~\ref{rem:dbeta:nearest}) replace the whole of the source's section on $\xi$, including its theorem on the threshold $\alpha_1=7.7892$, its numerical tables for $\alpha\in(1,\alpha_1)$ and its first open problem, which Proposition~\ref{bp:prop:valpha} settles. The equivalence of RH with positivity of the uncompensated density for all $\alpha>\frac12$ is Corollary~\ref{pol:cor:equiv} and Corollary~\ref{pol:cor:qrhggc} (Chapter~\ref{chap:polyazeros}), and Lagarias's criterion is cited from there. The $2\sin^2$ lemma is \eqref{eq:sin2} and Lemma~\ref{rz:lem:P1}, and its Thorin consequence for a L\'evy measure is \eqref{gl:eq:sinsq-thorin} and Theorem~\ref{gl:thm:sinsq}; the de Bruijn Newman theorem is Theorem~\ref{rz:thm:dbn}. The kink mechanism for weighted theta sums is that of Chapter~\ref{chap:miyatake} (Lemma~\ref{mi:lem:kink}, Theorem~\ref{mi:thm:kink}, Proposition~\ref{mi:prop:kinkbudget}, Proposition~\ref{mi:prop:reweight}). What is not elsewhere in the book is the identity decoded as Spitzer's formula together with the functional equation, the exact expansion of the shifted Comtet Majumdar family with its obstructions, the Gauss weights in the kink theorem with the explicit coefficients, the positivity of the Thorin density of the uniform law for every tilt, its Mellin transform, and the open problems.

Status words follow the rest of the book. A statement is proved when a proof or a pointer is given, computed when a script of the book produced it, conditional, heuristic or open as stated. The results are these.
\begin{enumerate}
\item The identity of Comtet and Majumdar holds on the whole critical strip, $0<\operatorname{Re}s<1$, and is the functional equation of $\zeta$ combined with Spitzer's formula for the maximum of a symmetric stable walk (Proposition~\ref{cg:prop:mellin}, Proposition~\ref{cg:prop:spitzer}, proved). The logarithm of the Laplace transform of the uniform law has a regularised signed L\'evy representation whose $k$ measure is Lebesgue measure minus counting measure, and $\zeta$ is its Mellin transform (Proposition~\ref{cg:prop:signed}, proved as improper integrals with a common cutoff).
\item An even entire function of order less than two has a GGC reciprocal exactly when its zeros are real (Lemma~\ref{cg:lem:lpggc}), so a family with GGC reciprocals converging to $\Xi$ proves RH (Corollary~\ref{cg:cor:closure}, proved).
\item The shifted Comtet Majumdar family $\zeta_\delta$ has an exact expansion in $\zeta(s+k)$ (Theorem~\ref{cg:thm:expansion}, proved). Its completion has a pole at $s=3$ for every $\delta\in(0,2\pi)$, so its reciprocal is not a GGC transform and not the transform of any positive measure (Proposition~\ref{cg:prop:obstructions}, proved; the claim of poles at all odd integers of the source is false and is repaired). Each simple zero moves at an explicit first order rate (Proposition~\ref{cg:prop:obstructions}(b), proved), and the first three leave the critical line (computed, Table~\ref{cg:tab:drift}). The shifted family is not the literal cutoff of Comtet and Majumdar, and the two differ already in sign (Remark~\ref{cg:rem:literal}, proved).
\item The uniform law, tilted by $\delta>0$ in the sense of the $2\sin^2$ lemma, has a strictly positive Thorin density for every $\delta>0$, found by Poisson summation and convexity of the Cauchy kernel (Theorem~\ref{cg:thm:uniform}, proved; the numerical constant of the proof is computed), and the regularised zeta function is its Mellin transform on $1<\operatorname{Re}s<3$ (Theorem~\ref{cg:thm:mellinthorin}, proved).
\item Gauss weights on the terms of the theta series produce, for every $\delta>0$, a kernel transform with finitely many real zeros and infinitely many nonreal zeros (Theorem~\ref{cg:thm:kink}, proved); symmetrising instead produces poles (Section~\ref{cg:sec:kernel}).
\item The de Bruijn Newman family is a deformation whose GGC property is Theorem~\ref{rz:thm:dbn}; what is false in the source, a convergence criterion for RH and the normalisation of the thresholds, is repaired (Section~\ref{cg:sec:dbn}). The remaining problems are open.
\end{enumerate}
